\section{Operating-Condition Reference Experiments}
\label{app:condition}

Section~\ref{sec:related} states that the leave-one-condition-out protocol,
the standard robustness test in the multi-condition literature, is nearly
saturated on this dataset, whereas the reverse direction is not. This
appendix reports the experiments behind that statement. They use a
deliberately generic pipeline---distinct from, and simpler than, the
constituent-level methodology of Section~\ref{sec:methods}---so that the
measured gaps reflect protocol difficulty rather than modeling choices, and
they are included with the released code.

Windows of 0.64~s (8192 samples, non-overlapping) were described by 108
time- and frequency-domain features (eight temporal statistics and ten
spectral descriptors per channel over the three vibration and three current
channels), classified into the original 24 diagnostic states by a 300-tree
random forest, with three seeds per configuration. Six split protocols were
compared. \emph{Random windows} splits windows irrespective of recording (a
deliberately leakage-prone reference); \emph{in-condition} trains on the
first 60\% of every recording and tests on the final 30\% with a 10\% guard
gap; \emph{unknown condition} holds out each of the twelve operating
conditions in turn; \emph{cross-profile} trains on one excitation profile
(constant-torque variable-speed or constant-speed variable-torque) and
tests on the other; \emph{steady~$\rightarrow$~transitional} trains on
quasi-stationary segments and tests on speed/load ramps, with stationarity
labeled from the key-phase-derived speed and the torque channel (the torque
channel is used only for this protocol construction and never as a
classifier input); and \emph{single source} inverts the unknown-condition
protocol, training on one condition and testing on the remaining eleven.

\begin{table}[!t]
\caption{Reference protocols, 24-class accuracy (mean $\pm$ std over folds
and three seeds). The same features and classifier are used throughout;
only the split changes.}
\label{tab:condref}
\centering
\footnotesize
\setlength{\tabcolsep}{5pt}
\begin{tabular}{l c c}
\toprule
Protocol & Folds & Accuracy \\
\midrule
Random windows (leakage-prone reference)   & 1  & $0.975 \pm 0.000$ \\
In-condition (guard-gap temporal split)    & 1  & $0.965 \pm 0.000$ \\
Unknown condition (leave-one-out)          & 12 & $0.939 \pm 0.024$ \\
Cross-profile                              & 2  & $0.839 \pm 0.115$ \\
Steady $\rightarrow$ transitional segments & 1  & $0.766 \pm 0.001$ \\
Single source (train on one condition)     & 12 & $0.333 \pm 0.104$ \\
\bottomrule
\end{tabular}
\end{table}

Table~\ref{tab:condref} summarizes the outcome. Holding one condition out
of twelve costs only ${\sim}2.6$ accuracy points relative to the
in-condition reference, because the eleven retained conditions bracket the
held-out one in speed and load and the model interpolates. Training on a
single condition and extrapolating to the other eleven costs over 60
points; paired over the twelve condition folds, the difference between the
two directions is 0.606 (paired $t$-test $p = 7.7\times10^{-10}$; Wilcoxon
$p = 4.9\times10^{-4}$). Transfer between excitation profiles is likewise
asymmetric at equal training size: training under variable speed
generalizes to variable load (accuracy 0.944), whereas the reverse loses
${\sim}21$ points (0.734; paired over seeds, $p = 3.3\times10^{-5}$).
Finally, appending 36 speed-invariant envelope-order descriptors to the
feature set leaves the interpolating protocols unchanged ($+0.001$ on
unknown-condition, paired $p = 0.84$) but significantly improves the
extrapolating one ($+0.047$ on single source, paired
$p = 8.4\times10^{-4}$), consistent with their construction.

Two implications carry into the main study. Methodologically, robustness to
operating conditions should be claimed in the extrapolating direction,
which motivates crossing the composition and context regimes with
unseen-condition evaluation rather than relying on leave-one-condition-out
alone. Practically, excitation diversity during commissioning matters more
than data volume: a training condition in which speed varies transfers to
load variation, but not conversely.
